% ============================================================
%  RAPPORT TECHNIQUE N°1 - Deep Learning
%  Architecture des Réseaux de Neurones & Problème XOR
%  Auteurs : Souaibou DIENG & Ibrahima FALL
%  ESP DAKAR - Master 2 SRT - 2025/2026
% ============================================================

\documentclass[12pt, a4paper]{article}

% ---------- Encodage & Langue ----------
\usepackage[utf8]{inputenc}
\usepackage[T1]{fontenc}
\usepackage[french]{babel}

% ---------- Mise en page ----------
\usepackage[top=2.5cm, bottom=2.5cm, left=3cm, right=2.5cm]{geometry}
\usepackage{setspace}
\onehalfspacing

% ---------- Polices ----------
\usepackage{lmodern}
\usepackage{microtype}

% ---------- Graphiques & Couleurs ----------
\usepackage{graphicx}
\usepackage{xcolor}
\usepackage{float}
\usepackage{tikz}
\usetikzlibrary{positioning, arrows.meta, shapes.geometric}

% Palette de couleurs (identique au rapport Fine-Tuning)
\definecolor{espblue}{RGB}{0, 70, 127}
\definecolor{espgold}{RGB}{204, 153, 0}
\definecolor{codegray}{RGB}{40, 44, 52}
\definecolor{codegreen}{RGB}{80, 200, 120}
\definecolor{codeorange}{RGB}{230, 126, 34}
\definecolor{codecomment}{RGB}{127, 140, 141}
\definecolor{backcolour}{RGB}{242, 244, 246}

% ---------- Tableaux ----------
\usepackage{booktabs}
\usepackage{tabularx}
\usepackage{array}
\usepackage{colortbl}

% ---------- Maths ----------
\usepackage{amsmath}
\usepackage{amssymb}
\usepackage{amsfonts}

% ---------- Code source ----------
\usepackage{listings}
\lstdefinestyle{pythonstyle}{
    backgroundcolor=\color{codegray},
    basicstyle=\ttfamily\footnotesize\color{white},
    keywordstyle=\color{codeorange}\bfseries,
    stringstyle=\color{codegreen},
    commentstyle=\color{codecomment}\itshape,
    numberstyle=\tiny\color{gray},
    numbers=left,
    numbersep=8pt,
    breaklines=true,
    breakatwhitespace=true,
    frame=single,
    framesep=5pt,
    rulecolor=\color{espblue},
    showstringspaces=false,
    tabsize=4,
    captionpos=b,
    literate=
        {é}{{\'{e}}}1 {è}{{\`{e}}}1 {ê}{{\^{e}}}1 {ë}{{\"e}}1
        {à}{{\`{a}}}1 {â}{{\^{a}}}1 {î}{{\^i}}1 {ô}{{\^o}}1
        {ù}{{\`{u}}}1 {û}{{\^{u}}}1 {ç}{{\c{c}}}1
}
\lstset{style=pythonstyle}

% ---------- Hyperliens ----------
\usepackage[hidelinks]{hyperref}
\hypersetup{
    pdftitle={Rapport TP1 - Architecture Réseaux de Neurones & XOR},
    pdfauthor={Souaibou DIENG \& Ibrahima FALL},
    pdfsubject={Deep Learning - Master 2 SRT},
}

% ---------- En-têtes et pieds de page ----------
\usepackage{fancyhdr}
\pagestyle{fancy}
\fancyhf{}
\fancyhead[L]{\small\textcolor{espblue}{\textit{Architecture Réseaux de Neurones \& Problème XOR}}}
\fancyhead[R]{\small\textcolor{espblue}{\textit{Master 2 SRT -- ESP Dakar}}}
\fancyfoot[C]{\textcolor{espblue}{\thepage}}
\renewcommand{\headrulewidth}{0.4pt}
\renewcommand{\headrule}{\hbox to\headwidth{%
    \color{espblue}\leaders\hrule height \headrulewidth\hfill}}

% ---------- Boîtes ----------
\usepackage{tcolorbox}
\tcbuselibrary{skins, breakable}

\newtcolorbox{infobox}[1][]{
    enhanced, breakable,
    colback=espblue!5,
    colframe=espblue,
    boxrule=0.8pt,
    arc=4pt,
    fonttitle=\bfseries\color{white},
    colbacktitle=espblue,
    #1
}

\newtcolorbox{mathbox}{
    enhanced,
    colback=espgold!8,
    colframe=espgold,
    boxrule=0.8pt,
    arc=4pt,
}

\newtcolorbox{resultbox}{
    enhanced,
    colback=codegreen!8,
    colframe=codegreen!60!black,
    boxrule=0.8pt,
    arc=4pt,
}

\newtcolorbox{casebox}[1]{
    enhanced,
    colback=espblue!4,
    colframe=espblue!40,
    boxrule=0.6pt,
    arc=3pt,
    title={\textbf{Cas #1}},
    fonttitle=\small\bfseries\color{espblue},
    colbacktitle=espblue!15,
    coltitle=espblue,
}

% ---------- Sections personnalisées ----------
\usepackage{titlesec}
\titleformat{\section}
    {\normalfont\Large\bfseries\color{espblue}}
    {\thesection.}{10pt}{}
    [\vspace{2pt}{\color{espblue}\titlerule[0.6pt]}]
\titlespacing{\section}{0pt}{18pt}{10pt}

\titleformat{\subsection}
    {\normalfont\large\bfseries\color{espblue!80!black}}
    {\thesubsection}{8pt}{}
\titlespacing{\subsection}{0pt}{12pt}{6pt}

\titleformat{\paragraph}[runin]
    {\normalfont\normalsize\bfseries\color{espblue!70!black}}
    {}{0pt}{}[\ ---]

% ---------- Divers ----------
\usepackage{enumitem}
\usepackage{caption}
\captionsetup{font=small, labelfont=bf, labelsep=period}

% ============================================================
%  DÉBUT DU DOCUMENT
% ============================================================
\begin{document}

% ============================================================
%  PAGE DE GARDE
% ============================================================
\begin{titlepage}
    \centering
    \vspace*{0.3cm}

    {\color{espblue}\rule{\linewidth}{3pt}}\\[0.5cm]

    % Entête institutionnel avec logo centré
    \begin{minipage}[c]{0.30\textwidth}
        \centering
        {\small\bfseries RÉPUBLIQUE DU SÉNÉGAL}\\[3pt]
        {\footnotesize Un Peuple – Un But – Une Foi}
    \end{minipage}
    \hfill
    \begin{minipage}[c]{0.30\textwidth}
        \centering
        \includegraphics[width=2.8cm]{logo_esp.png}
    \end{minipage}
    \hfill
    \begin{minipage}[c]{0.30\textwidth}
        \centering
        {\small\bfseries UNIVERSITÉ CHEIKH ANTA DIOP}\\[3pt]
        {\footnotesize DE DAKAR (UCAD)}
    \end{minipage}

    \vspace{0.5cm}
    {\color{espblue}\rule{\linewidth}{1pt}}\\[0.3cm]

    {\large\bfseries\color{espblue} ÉCOLE SUPÉRIEURE POLYTECHNIQUE DE DAKAR}\\[4pt]
    {\normalsize DÉPARTEMENT DE GÉNIE INFORMATIQUE}\\[2pt]
    {\normalsize\bfseries MASTER 2 -- SYSTÈMES, RÉSEAUX ET TÉLÉCOMMUNICATIONS}

    \vspace{0.4cm}
    {\color{espgold}\rule{0.5\linewidth}{2pt}}
    \vspace{0.3cm}

    {\large\color{espblue!70!black} Module :}\\[3pt]
    {\Large\bfseries Deep Learning / Traitement d'Images}

    \vspace{0.3cm}
    {\color{espgold}\rule{0.8\linewidth}{1.5pt}}
    \vspace{0.3cm}

    % Badge numéro de TP
    \begin{tcolorbox}[enhanced, width=4cm, colback=espblue, colframe=espblue,
        arc=20pt, boxrule=0pt, halign=center]
        {\large\bfseries\color{white} Rapport Technique}\\[2pt]
        {\LARGE\bfseries\color{espgold} N° 1}
    \end{tcolorbox}

    \vspace{0.3cm}

    {\LARGE\bfseries\color{espblue}
        Architecture des Réseaux de Neurones\\[5pt]
        et Problème du XOR
    }

    \vspace{0.3cm}
    {\color{espgold}\rule{0.8\linewidth}{1.5pt}}
    \vspace{0.5cm}

    \Large\bfseries Réalisé par :\\[0.4cm]
    \begin{tabular}{cc}
        \Large\color{espblue} Souaibou \textsc{Dieng} &
        \Large\color{espblue} Ibrahima \textsc{Fall} \\
    \end{tabular}

    \vfill

    \begin{tabular}{rl}
        \textbf{Classe :}              & Master 2 SRT \\[3pt]
        \textbf{Encadrant :}           & Dr. Mouhamed Kaba \textsc{Keïta} \\[3pt]
        \textbf{Année universitaire :} & 2025 -- 2026 \\
    \end{tabular}

    \vspace{0.4cm}
    {\color{espblue}\rule{\linewidth}{3pt}}
\end{titlepage}

% ============================================================
%  TABLE DES MATIÈRES
% ============================================================
\tableofcontents
\thispagestyle{fancy}
\newpage

% ============================================================
%  INTRODUCTION
% ============================================================
\section*{Introduction}
\addcontentsline{toc}{section}{Introduction}

L'apprentissage profond repose sur la capacité des neurones artificiels à modéliser
des fonctions complexes. Cependant, les modèles simples comme le perceptron de Rosenblatt
présentent des limites théoriques majeures, notamment l'impossibilité de résoudre des
problèmes qui ne sont pas \textbf{linéairement séparables}.

Ce rapport traite des fondamentaux du Deep Learning à travers l'étude du problème logique
\textbf{XOR} (\textit{OU exclusif}). Nous détaillerons l'architecture d'un
\textbf{Perceptron Multicouche (MLP)}, les calculs de propagation avant
(\textit{forward propagation}), et nous proposerons une implémentation logicielle
permettant de valider la théorie.

% ============================================================
%  SECTION 1 : MODÉLISATION
% ============================================================
\section{Modélisation et Séparabilité Linéaire}

\subsection{Le Neurone Artificiel}

Un neurone artificiel effectue une combinaison linéaire de ses entrées pondérées,
suivie de l'application d'une \textbf{fonction d'activation} $f$ :

\begin{mathbox}
\begin{equation}
    y = f\!\left(\sum_{i=1}^{n} w_i\, x_i + b\right)
    \label{eq:neuron}
\end{equation}
où $x_i$ sont les entrées, $w_i$ les poids synaptiques, $b$ le biais et $f$ la
fonction d'activation.
\end{mathbox}

\subsection{Limites du Perceptron Simple : le cas du XOR}

Le XOR est une fonction logique dont la table de vérité est la suivante :

\begin{table}[H]
    \centering
    \caption{Table de vérité du XOR}
    \label{tab:xor}
    \renewcommand{\arraystretch}{1.4}
    \begin{tabular}{ccc}
        \toprule
        \rowcolor{espblue!15}
        \textbf{Entrée $x_1$} & \textbf{Entrée $x_2$} & \textbf{Sortie attendue} \\
        \midrule
        0 & 0 & 0 \\
        0 & 1 & \cellcolor{codegreen!20}\textbf{1} \\
        1 & 0 & \cellcolor{codegreen!20}\textbf{1} \\
        1 & 1 & 0 \\
        \bottomrule
    \end{tabular}
\end{table}

\begin{infobox}[title={Pourquoi un perceptron simple ne suffit pas ?}]
Contrairement aux fonctions AND ou OR, il est \textbf{impossible} de tracer une seule
droite (hyperplan) séparant les points de sortie 0 des points de sortie 1 dans l'espace
$\mathbb{R}^2$. Cette \textbf{non-linéarité} impose l'utilisation d'au moins une
\textbf{couche cachée}.
\end{infobox}

% ============================================================
%  SECTION 2 : ARCHITECTURE MLP
% ============================================================
\section{Architecture du Perceptron Multicouche}

Pour résoudre le XOR, nous utilisons un réseau composé de :
\begin{itemize}[leftmargin=1.8cm]
    \item une \textbf{couche d'entrée} : 2 neurones ($x_1$, $x_2$)
    \item une \textbf{couche cachée} : 2 neurones ($N_1$, $N_2$)
    \item une \textbf{couche de sortie} : 1 neurone ($y$)
\end{itemize}

\begin{figure}[H]
    \centering
    \includegraphics[width=0.82\textwidth]{arch_tp1.png}
    \caption{Architecture du MLP proposée pour le XOR avec poids et biais configurés.}
    \label{fig:xor_arch}
\end{figure}

\subsection{Paramétrage du Réseau}

Selon l'architecture de la Figure~\ref{fig:xor_arch}, les paramètres sont fixés comme suit :

\begin{table}[H]
    \centering
    \caption{Paramètres du réseau MLP pour le XOR}
    \label{tab:params}
    \renewcommand{\arraystretch}{1.5}
    \begin{tabularx}{0.85\textwidth}{lccX}
        \toprule
        \rowcolor{espblue!15}
        \textbf{Neurone} & \textbf{Poids $w$} & \textbf{Biais $b$} & \textbf{Fonction} \\
        \midrule
        $N_1$ (couche cachée)  & $[1,\; 1]$   & $-1.5$ & Heaviside \\
        $N_2$ (couche cachée)  & $[1,\; 1]$   & $-0.5$ & Heaviside \\
        $y$\; (sortie)         & $[-1,\; 1]$  & $-0.5$ & Heaviside \\
        \bottomrule
    \end{tabularx}
\end{table}

Nous utilisons la \textbf{fonction à seuil de Heaviside} comme fonction d'activation :

\begin{mathbox}
\begin{equation}
    f(z) = \begin{cases} 1 & \text{si } z \geq 0 \\ 0 & \text{sinon} \end{cases}
    \label{eq:heaviside}
\end{equation}
\end{mathbox}

% ============================================================
%  SECTION 3 : VÉRIFICATION MATHÉMATIQUE
% ============================================================
\section{Vérification des Cas d'Usage}

Démontrons par le calcul que cette configuration résout effectivement le XOR
pour les quatre combinaisons d'entrées possibles.

% --- Cas (0,0) ---
\begin{casebox}{$(x_1, x_2) = (0, 0)$ — Sortie attendue : \textbf{0}}
\begin{align*}
    N_1 &= f\!\left(0 \cdot 1 + 0 \cdot 1 - 1.5\right) = f(-1.5) = \mathbf{0} \\
    N_2 &= f\!\left(0 \cdot 1 + 0 \cdot 1 - 0.5\right) = f(-0.5) = \mathbf{0} \\
    y   &= f\!\left(0 \cdot (-1) + 0 \cdot 1 - 0.5\right) = f(-0.5) = \mathbf{0} \checkmark
\end{align*}
\end{casebox}

\vspace{0.5cm}

% --- Cas (1,0) ---
\begin{casebox}{$(x_1, x_2) = (1, 0)$ — Sortie attendue : \textbf{1}}
\begin{align*}
    N_1 &= f\!\left(1 \cdot 1 + 0 \cdot 1 - 1.5\right) = f(-0.5) = \mathbf{0} \\
    N_2 &= f\!\left(1 \cdot 1 + 0 \cdot 1 - 0.5\right) = f(+0.5) = \mathbf{1} \\
    y   &= f\!\left(0 \cdot (-1) + 1 \cdot 1 - 0.5\right) = f(+0.5) = \mathbf{1} \checkmark
\end{align*}
\end{casebox}

\vspace{0.5cm}

% --- Cas (0,1) ---
\begin{casebox}{$(x_1, x_2) = (0, 1)$ — Sortie attendue : \textbf{1}}
\begin{align*}
    N_1 &= f\!\left(0 \cdot 1 + 1 \cdot 1 - 1.5\right) = f(-0.5) = \mathbf{0} \\
    N_2 &= f\!\left(0 \cdot 1 + 1 \cdot 1 - 0.5\right) = f(+0.5) = \mathbf{1} \\
    y   &= f\!\left(0 \cdot (-1) + 1 \cdot 1 - 0.5\right) = f(+0.5) = \mathbf{1} \checkmark
\end{align*}
\end{casebox}

\vspace{0.5cm}

% --- Cas (1,1) ---
\begin{casebox}{$(x_1, x_2) = (1, 1)$ — Sortie attendue : \textbf{0}}
\begin{align*}
    N_1 &= f\!\left(1 \cdot 1 + 1 \cdot 1 - 1.5\right) = f(+0.5) = \mathbf{1} \\
    N_2 &= f\!\left(1 \cdot 1 + 1 \cdot 1 - 0.5\right) = f(+1.5) = \mathbf{1} \\
    y   &= f\!\left(1 \cdot (-1) + 1 \cdot 1 - 0.5\right) = f(-0.5) = \mathbf{0} \checkmark
\end{align*}
\end{casebox}

\vspace{0.6cm}
\begin{resultbox}
\textbf{Résultat :} Les quatre cas de la table de vérité du XOR sont correctement
résolus par ce réseau MLP à deux neurones cachés.
\end{resultbox}

% ============================================================
%  SECTION 4 : IMPLÉMENTATION PYTHON
% ============================================================
\section{Implémentation Logicielle}

Le code suivant implémente la logique du réseau XOR en Python avec la bibliothèque
\texttt{numpy}, en respectant exactement les poids et biais définis à la
Section~\ref{tab:params}.

\begin{lstlisting}[language=Python, caption={Simulateur du réseau XOR -- Python/NumPy}, label=lst:xor]
import numpy as np

def step_function(x):
    """Fonction de Heaviside : retourne 1 si x >= 0, sinon 0."""
    return 1 if x >= 0 else 0

def xor_network(x1, x2):
    # --- Couche cachee ---
    n1 = step_function(x1 * 1 + x2 * 1 - 1.5)  # N1 : AND approxime
    n2 = step_function(x1 * 1 + x2 * 1 - 0.5)  # N2 : OR

    # --- Couche de sortie ---
    y = step_function(n1 * (-1) + n2 * 1 - 0.5) # XOR = OR AND NOT(AND)
    return y

# Verification sur tous les cas
inputs = [(0, 0), (0, 1), (1, 0), (1, 1)]
expected = [0, 1, 1, 0]

print(f"{'X1':>4} {'X2':>4} | {'Attendu':>8} {'Obtenu':>8} {'OK':>4}")
print("-" * 36)
for (x1, x2), exp in zip(inputs, expected):
    res = xor_network(x1, x2)
    ok  = "V" if res == exp else "X"
    print(f"{x1:>4} {x2:>4} | {exp:>8} {res:>8} {ok:>4}")
\end{lstlisting}

La sortie attendue du programme est la suivante :

\begin{table}[H]
    \centering
    \caption{Résultats d'exécution du simulateur Python}
    \label{tab:results}
    \renewcommand{\arraystretch}{1.4}
    \begin{tabular}{ccccc}
        \toprule
        \rowcolor{espblue!15}
        \textbf{$x_1$} & \textbf{$x_2$} & \textbf{Attendu} & \textbf{Obtenu} & \textbf{Statut} \\
        \midrule
        0 & 0 & 0 & 0 & \textcolor{codegreen!60!black}{\textbf{OK}} \\
        0 & 1 & 1 & 1 & \textcolor{codegreen!60!black}{\textbf{OK}} \\
        1 & 0 & 1 & 1 & \textcolor{codegreen!60!black}{\textbf{OK}} \\
        1 & 1 & 0 & 0 & \textcolor{codegreen!60!black}{\textbf{OK}} \\
        \bottomrule
    \end{tabular}
\end{table}

% ============================================================
%  CONCLUSION
% ============================================================
\section*{Conclusion}
\addcontentsline{toc}{section}{Conclusion}

L'étude du problème XOR démontre qu'une architecture à plusieurs couches est
\textbf{indispensable} pour traiter des données non linéairement séparables.
En superposant des neurones simples, le réseau est capable de créer des
\textbf{frontières de décision complexes}.

\begin{resultbox}
Ce principe constitue la base du Deep Learning moderne : la capacité de
\textbf{transformer et projeter les données} dans de nouveaux espaces représentatifs
où la classification linéaire devient possible. C'est cette propriété que les
architectures profondes (CNN, Transformers) exploitent à très grande échelle.
\end{resultbox}

\end{document}